\documentclass[11pt]{amsart}
\usepackage{amssymb,mathtools,array}
\usepackage[margin=1.15in]{geometry}
\usepackage[hidelinks]{hyperref}
\newtheorem{theorem}{Theorem}[section]
\newtheorem{proposition}[theorem]{Proposition}
\newtheorem{lemma}[theorem]{Lemma}
\newtheorem{corollary}[theorem]{Corollary}
\theoremstyle{definition}
\newtheorem{definition}[theorem]{Definition}
\theoremstyle{remark}
\newtheorem{remark}[theorem]{Remark}
\newcommand{\UU}{\mathcal U}
\newcommand{\Set}{\mathrm{Set}}
\newcommand{\Hom}{\operatorname{Hom}}
\newcommand{\Iso}{\operatorname{Iso}}
\newcommand{\Fun}{\operatorname{Fun}}
\newcommand{\Nat}{\operatorname{Nat}}
\newcommand{\Alg}{\operatorname{Alg}}
\newcommand{\Ob}{\operatorname{Ob}}
\newcommand{\FP}{\operatorname{FP}}
\newcommand{\id}{\operatorname{id}}
\newcommand{\fib}{\operatorname{fib}}
\newcommand{\tr}[1]{\left\Vert #1\right\Vert}
\newcommand{\code}[1]{\texttt{\small #1}}
\title[Univalent categories and functors of points]{Small univalent categories and functors of points\\
Replacement, finite presentations, and qcqs schemes}
\date{October 29, 2026}
\begin{document}
\raggedbottom
\begin{abstract}
A category can be presented as a 1-type with sets of arrows, identities,
and composition, with its identity types canonically equivalent to its
sets of isomorphisms. We show how replacement produces the small full
category generated by a small family in a locally small category. Its
universal property follows from descent along fully faithful, essentially
surjective functors. Applied to finite presentations, this constructs a
small univalent category whose objects are precisely the finitely
presented commutative algebras over a fixed ring. Yoneda then gives the
fully faithful contravariant embedding into set-valued functors. We then
prove that full faithfulness and mere essential surjectivity between
univalent categories give an explicit categorical equivalence without
choice. Applied to qcqs schemes, this identifies the geometric category
of locally ringed lattices with the category of functors satisfying the
Zariski sheaf and finite affine atlas conditions, both absolutely and
over an arbitrary small base ring. The last section
records the correspondence with the checked implementation.
\end{abstract}
\maketitle

\section{Categories as structured 1-types}

We work in univalent type theory with a hierarchy of universes $\UU_\ell$.
A proposition is a type any two of whose elements are equal; a set is a
type whose identity types are propositions; a 1-type is a type whose
identity types are sets. We write $\tr{T}$ for propositional truncation,
so that $\tr{T}$ expresses mere existence of an element of $T$.
Function extensionality follows from univalence.

A \emph{precategory} $C$ consists of a type $\Ob(C)$, sets $C(x,y)$,
identities $1_x$, and composition
\[
 C(y,z)\times C(x,y)\longrightarrow C(x,z),\qquad (g,f)\longmapsto g\circ f,
\]
satisfying the two unit laws and associativity. Its objects need not form
a set. For an arrow $f$, the type of inverse arrows together with the two
inverse equations is a proposition: the inverse is unique, and the
inverse equations are propositions. Consequently $\Iso_C(x,y)$ is a set.

Path induction gives a canonical function
\[
 \operatorname{idtoiso}_{x,y}:(x=y)\longrightarrow\Iso_C(x,y),
 \qquad \operatorname{idtoiso}(\operatorname{refl}_x)=1_x.
\]
\begin{definition}
A \emph{category} is a precategory for which every function
$\operatorname{idtoiso}_{x,y}$ is an equivalence. Its inverse is denoted
$\operatorname{isotoid}$.
\end{definition}
It follows that the objects of a category form a 1-type. Thus we can
regard a category as a 1-type equipped with hom sets and category
operations, subject also to the requirement that the canonical maps
$(x=y)\to\Iso_C(x,y)$ be equivalences. The 1-type condition alone is
insufficient. All arrows remain part of the structure, including the
noninvertible ones. The equality $(x=x)\simeq\operatorname{Aut}_C(x)$
explains why this object type is not generally a set of isomorphism classes.

\pagebreak[3]
The following criterion lets us prove univalence by a uniqueness argument.
\begin{lemma}[Total-space criterion]\label{lem:total}
A precategory $C$ is a category if and only if, for every $x$, the type
\[
 S_x=\sum_{y:\Ob(C)}\Iso_C(x,y)
\]
is a proposition. In that case $S_x$ is contractible, with centre $(x,1_x)$.
\end{lemma}
\begin{proof}
In a category, $S_x$ is equivalent to $\sum_y(x=y)$, which is contractible.
Conversely, suppose $S_x$ is a proposition. For $e:x\cong y$, project the
unique equality $(x,1_x)=(y,e)$ onto the object coordinate to obtain
$p_e:x=y$. Path induction on this equality in $S_x$ gives
$\operatorname{idtoiso}(p_e)=e$. Path induction on $p:x=y$ gives
$p_{\operatorname{idtoiso}(p)}=p$: in the reflexive case the equality
$(x,1_x)=(x,1_x)$ is reflexive since $S_x$ is a proposition.
These are the required inverse laws.
\end{proof}

A functor is \emph{fully faithful} if each of its functions on hom sets
is an equivalence. It is \emph{essentially surjective} if
\[
 \prod_{b:\Ob(B)}\tr{\sum_{c:\Ob(C)}\Iso_B(qc,b)}
\]
is inhabited. We distinguish this from \emph{split} essential
surjectivity, which supplies an actual pair $(c,e)$ for each $b$.
Full faithfulness reflects isomorphisms: lift an inverse by fullness and
check its two equations by faithfulness.

\section{Functor categories}

For precategories $B,E$, functors $B\to E$ and natural transformations
form a precategory $\Fun(B,E)$. A natural transformation $\alpha:F\to G$
is a family of arrows $\alpha_b:Fb\to Gb$ such that
\[
 G(f)\circ\alpha_b=\alpha_{b'}\circ F(f)\qquad(f:b\to b').
\]
Naturality is a proposition. Transformations form a set, and are equal
when their components are equal. Their identities and composition are
componentwise.

\begin{proposition}\label{prop:fununiv}
If $E$ is a category and $B$ is any precategory, then $\Fun(B,E)$ is a
category. In particular,
\[
 (F=G)\simeq\Iso_{\Fun(B,E)}(F,G).
\]
\end{proposition}
\begin{proof}
An isomorphism of functors has isomorphisms as its components. Conversely,
if all components of a natural transformation are isomorphisms, their
inverses are natural by multiplying the naturality square by the two
inverses. They give its inverse transformation.

Fix $F$. Consider the total type of pairs $(G,\alpha)$ where $\alpha:F\cong G$.
Its object data are a family $G_0(b)$ and isomorphisms
$a_b:Fb\cong G_0(b)$. The type of such data is equivalent to
\[
 \prod_{b:\Ob(B)}\left(\sum_{y:\Ob(E)}\Iso_E(Fb,y)\right),
\]
which is contractible by Lemma~\ref{lem:total}. For fixed object data,
naturality forces the arrow action to be
\[
 G_1(f)=a_{b'}\circ F(f)\circ a_b^{-1}.
\]
Thus there is at most one arrow action. The functor laws and naturality
are propositions. The total type $(G,\alpha)$ is therefore a proposition,
and Lemma~\ref{lem:total} applies to $\Fun(B,E)$.
\end{proof}

\section{Replacement and the small full image}

A type is \emph{$\UU_\ell$-locally small} if each of its identity types is
equivalent to a type in $\UU_\ell$. We use the following form of
replacement. For $P:\UU_\ell$, a locally small type $T:\UU_k$, and a
function $F:P\to T$, there are
\[
 P\xrightarrow{q}I\xrightarrow{i}T,\qquad I:\UU_\ell,
 \qquad i(qp)=F(p),
\]
where $i$ is an embedding and $q$ is surjective. Here being an embedding
means that each function
\[
 (a=b)\longrightarrow(i(a)=i(b))
\]
is an equivalence. Surjectivity permits elimination into dependent
families of propositions in any universe. The small image $I$ is the
same for all these elimination universes. This last clause matters when
membership in the image is a proposition in a larger universe.

\begin{theorem}[Small full image]\label{thm:image}
Let $D$ be a category with hom sets in $\UU_\ell$, and let
$F:P\to\Ob(D)$ with $P:\UU_\ell$. There is a category $D_F$ with both
objects and hom sets in $\UU_\ell$, together with a fully faithful
inclusion $i:D_F\to D$, such that
\begin{equation}\label{eq:member}
 \fib_i(x)\simeq
 \tr{\sum_{p:P}\Iso_D(Fp,x)}.
\end{equation}
The fibre on the left is a proposition.
\end{theorem}
\begin{proof}
The object type of $D$ is locally small because
$(x=y)\simeq\Iso_D(x,y)$ and the latter type is small. Apply replacement
to $F$, obtaining $P\to I\to\Ob(D)$. Define
\[
 \Ob(D_F)=I,\qquad D_F(a,b)=D(ia,ib),
\]
with identities and composition inherited from $D$. The inclusion is the
identity on hom sets and is fully faithful. Its univalence follows from
\[
 (a=b)\simeq(ia=ib)\simeq\Iso_D(ia,ib)
                  =\Iso_{D_F}(a,b).
\]
By path induction this composite is the canonical $\operatorname{idtoiso}$.

Since $i$ is an embedding, $\fib_i(x)=\sum_{a:I}(ia=x)$ is a proposition.
A pair $(p,e:Fp\cong x)$ determines an element of this fibre by univalence
and $i(qp)=Fp$. We can therefore eliminate its mere existence into the
fibre. Conversely, a fibre point $(a,r)$ and surjectivity of $q$ give
merely some $p$ with $qp=a$, hence an isomorphism $Fp\cong x$.
Both sides of \eqref{eq:member} are propositions, so these maps give an
equivalence.
\end{proof}

There is also a precategory $C_F$ with objects $P$ and arrows
$C_F(p,p')=D(Fp,Fp')$. The map $q$ extends to a functor
$q:C_F\to D_F$. More explicitly, write
$t_p:i(qp)\cong Fp$ for the isomorphism given by the image equation. Then
\[
 q(h)=t_{p'}^{-1}\circ h\circ t_p.
\]
Conjugation is an equivalence on hom sets, and every object of $D_F$ is
merely isomorphic to some $qp$. Thus $q$ is fully faithful and essentially
surjective. We next give its universal property.

\section{Descent along a fully faithful, essentially surjective functor}

Throughout this section let $q:C\to B$ be fully faithful and essentially
surjective. The source and target may initially be precategories.

\begin{lemma}\label{lem:ff}
For any precategory $E$, precomposition
\[
 q^*:\Fun(B,E)\longrightarrow\Fun(C,E)
\]
is fully faithful.
\end{lemma}
\begin{proof}
Let $G,H:B\to E$ and $\beta:Gq\to Hq$ be natural. For each $b$, seek an
arrow $\gamma:Gb\to Hb$ satisfying, for every $c$ and every arrow
$u:qc\to b$,
\begin{equation}\label{eq:extendnat}
 \gamma\circ G(u)=H(u)\circ\beta_c.
\end{equation}
The type $T_b$ of an arrow with these equations is a proposition.
Indeed, to compare two arrows in it, we may merely choose
$e:qc\cong b$, since their equality is a proposition. Equation
\eqref{eq:extendnat} at $e$ and cancellation of the isomorphism $G(e)$
give equality; the equations attached to the arrows are propositions.

To inhabit $T_b$, we can now eliminate essential surjectivity into this
proposition. For a given $e_0:qc_0\cong b$, put
\[
 \gamma_0=H(e_0)\circ\beta_{c_0}\circ G(e_0)^{-1}.
\]
For any $u:qc\to b$, full faithfulness gives the unique $f:c\to c_0$
with $qf=e_0^{-1}u$. Naturality of $\beta$ proves
\eqref{eq:extendnat} for $u=e_0\circ qf$. Hence $T_b$ is contractible.
Its unique arrow defines $\alpha_b$.

For $g:b\to b'$, compare $H(g)\alpha_b$ and $\alpha_{b'}G(g)$ after
composition with $G(e)$ for a merely chosen $e:qc\cong b$. Their equality
follows from \eqref{eq:extendnat} at $e$ and at $ge$. Cancellation gives
naturality. At $b=qc$, taking $u=1_{qc}$ gives $\alpha_{qc}=\beta_c$.
Finally, any natural transformation restricting to $\beta$ satisfies
\eqref{eq:extendnat}, so it equals $\alpha$. This constructs an inverse
to precomposition on each hom set.
\end{proof}

\begin{lemma}\label{lem:descent}
If $E$ is a category, every functor $G:C\to E$ has a descent
$H:B\to E$ with a natural isomorphism $Hq\cong G$.
\end{lemma}
\begin{proof}
We construct $Hb$ without choosing a representative of $b$. Define $X_b$
to consist of an object $x:E$ and isomorphisms
\[
 \eta_{c,e}:Gc\cong x\qquad(c:C,\ e:qc\cong b)
\]
satisfying
\begin{equation}\label{eq:cocone}
 \eta_{c',e'}\circ G(f)=\eta_{c,e}
 \quad\text{whenever } e'\circ qf=e.
\end{equation}
We first prove $X_b$ is a proposition. This assertion is itself a
proposition, so essential surjectivity lets us suppose
$e_0:qc_0\cong b$ is given. Evaluation at $(c_0,e_0)$ maps $X_b$ to
\[
 S_{Gc_0}=\sum_{x:E}\Iso_E(Gc_0,x).
\]
It has a reconstruction: from $(x,\varphi)$ define
\[
 \eta_{c,e}=\varphi\circ G(f_e),\qquad qf_e=e_0^{-1}\circ e.
\]
Full faithfulness gives $f_e$, and reflects its invertibility.
Faithfulness and the functor laws give \eqref{eq:cocone}. That same
equation shows that reconstructing an element of $X_b$ from its evaluation
gives the original element. Hence $X_b$ is a retract of $S_{Gc_0}$,
which is a proposition by Lemma~\ref{lem:total}. Thus $X_b$ is a proposition.

Only now do we use essential surjectivity to inhabit $X_b$: given
$(c_0,e_0)$, reconstruct from $(Gc_0,1)$. We obtain a unique element of
$X_b$, denoted $(Hb,\eta_b)$.

For $g:b\to b'$, define $Y_g$ to be the type of arrows
$\varphi:Hb\to Hb'$ satisfying
\begin{equation}\label{eq:descendarrow}
 \varphi\circ\eta_{b,c,e}
   =\eta_{b',c',e'}\circ G(f)
 \quad\text{if }e'\circ qf=g\circ e.
\end{equation}
This is a proposition. To compare two such arrows, merely choose
$e:qc\cong b$ and $e':qc'\cong b'$, lift $e'^{-1}ge$ to $f$, and cancel
$\eta_{b,c,e}$. To inhabit it, use the same representatives and define
\[
 \varphi_0=\eta_{b',c',e'}\circ G(f)\circ\eta_{b,c,e}^{-1}.
\]
For other representatives $(d,u),(d',u')$, lift $e^{-1}u$ and
$e'^{-1}u'$ to $m:d\to c$ and $m':d'\to c'$. If
$u'\circ qh=g\circ u$, faithfulness gives $fm=m'h$.
Equation \eqref{eq:cocone} and this equality verify
\eqref{eq:descendarrow} for $\varphi_0$. Thus $Y_g$ is contractible;
its arrow defines $H(g)$.

The identity satisfies \eqref{eq:descendarrow} by
\eqref{eq:cocone}, so $H(1)=1$. For composable $g,g'$, merely choose a
representative for their intermediate object. Lift the two arrows and
use faithfulness to identify the composite of the lifts with a lift of
$g'g$. Then $H(g')H(g)$ satisfies \eqref{eq:descendarrow} for $g'g$,
so uniqueness gives the composition law. Finally,
$\eta_{qc,c,1}:Gc\cong H(qc)$ is natural by
\eqref{eq:descendarrow}. Its inverse gives $Hq\cong G$.
\end{proof}

The order of these arguments is essential: the types of possible values
are proved to be propositions before a merely existing representative
is used to inhabit them. No choice of representatives is required.

\begin{proposition}\label{prop:equivobjects}
A fully faithful, split essentially surjective functor $P:X\to Y$
between categories induces an equivalence $\Ob(X)\simeq\Ob(Y)$.
\end{proposition}
\begin{proof}
For $y:Y$, the type $S_y=\sum_{x:X}\Iso_Y(Px,y)$ is a proposition.
Indeed, for $(x,e),(x',e')$, lift the isomorphism $e'^{-1}e$ along $P$
to $\varphi:x\cong x'$. Univalence of $X$ makes it a path $x=x'$.
Path induction on this path, using $e'P\varphi=e$, identifies the two
pairs. Univalence of $Y$ identifies $S_y$ with the fibre
$\sum_x(Px=y)$. The splitting inhabits that fibre, so every fibre is
contractible.
\end{proof}

\begin{theorem}[From mere essential surjectivity to equivalence]
\label{thm:weak-equivalence}
Let $P:X\to Y$ be a fully faithful, essentially surjective functor between
univalent categories. Then its object map is an equivalence of types,
and $P$ has a quasi-inverse functor with natural isomorphisms for both
composites. Split essential surjectivity is not an additional hypothesis.
\end{theorem}
\begin{proof}
Full faithfulness gives equivalences on isomorphism types: lift an
isomorphism and its inverse by fullness, and prove the inverse equations
by faithfulness. The square
\[
\begin{array}{ccc}
 (x=x')&\xrightarrow{\operatorname{ap}P}&(Px=Px')\\
 \downarrow\scriptstyle{\simeq}&&\downarrow\scriptstyle{\simeq}\\
 \Iso_X(x,x')&\xrightarrow{\simeq}&\Iso_Y(Px,Px')
\end{array}
\]
commutes by path induction. Its vertical maps are the canonical
univalence equivalences, so $\operatorname{ap}P$ is an equivalence.
Consequently the object map of $P$ is an embedding, and every fibre
\[
 T_y=\sum_{x:\Ob(X)}(Px=y)
\]
is a proposition. Essential surjectivity and univalence of $Y$ give
$\tr{T_y}$. Eliminate this truncation into the already propositional
$T_y$. Each fibre is now inhabited and contractible, which proves the
asserted equivalence on object types.

Write its fibre centres as $(Ky,p_y:PKy=y)$, and let
$e_y:PKy\cong y$ be the path-isomorphism of $p_y$. For $a:y\to z$,
define $Ka$ by the unique lift
\begin{equation}\label{eq:inverse-arrow}
 P(Ka)=e_z^{-1}\circ a\circ e_y.
\end{equation}
Equivalently, $Ka$ is the unique arrow satisfying
$e_z\circ P(Ka)=a\circ e_y$. The identity arrow satisfies this square
for $1_y$, and pasting the squares for two composable arrows proves the
square for their composite. Uniqueness proves the two functor laws.
The squares themselves state that $e$ is a natural isomorphism
$PK\cong1_Y$. Reflect $e_{Px}$ along the fully faithful $P$ to obtain
$KPx\cong x$. Applying faithfulness to the naturality of $e$ proves
naturality of these reflected isomorphisms, giving $KP\cong1_X$.
Thus $K$ is a quasi-inverse. Every truncation was eliminated into a
proposition; no axiom of choice is used.
\end{proof}

\begin{theorem}[Universal property of descent]\label{thm:descent}
For every category $E$, precomposition by $q$ is fully faithful and split
essentially surjective. It induces an equivalence of types
\[
 \Fun(B,E)\simeq\Fun(C,E).
\]
\end{theorem}
\begin{proof}
Lemmas~\ref{lem:ff} and \ref{lem:descent} give the two assertions about
precomposition. Both functor categories are categories by
Proposition~\ref{prop:fununiv}; apply Proposition~\ref{prop:equivobjects}.
\end{proof}

Applying this to $C_F\to D_F$ gives
\begin{equation}\label{eq:completion}
 \Fun(D_F,E)\simeq\Fun(C_F,E).
\end{equation}
Thus $q:C_F\to D_F$ has the universal property of the \emph{univalent
completion}, also called the \emph{Rezk completion}, of $C_F$.
Here the existing ambient category $D$ and replacement provide a small
construction of that completion. Replacement supplies the small object
type; the descent argument supplies the universal property.

\section{The category of commutative algebras}

Fix a commutative unital ring $R$ whose carrier is a set in $\UU_\ell$.
The zero ring is allowed. An $R$-algebra is a commutative unital ring $A$
with carrier a set in $\UU_\ell$, together with a unital ring map
$\iota_A:R\to A$. An $R$-algebra map $f:A\to B$ is a unital ring map
satisfying $f\iota_A=\iota_B$. Such maps form a set, denoted
$\Hom_R(A,B)$. Pointwise identity and composition define a precategory
$\Alg_R$ with objects in $\UU_{\ell+1}$ and hom sets in $\UU_\ell$.

\begin{lemma}\label{lem:algiso}
An $R$-algebra map is an isomorphism if and only if its underlying
function is an equivalence of sets.
\end{lemma}
\begin{proof}
An algebraic inverse is a function inverse. Conversely, let $g$ be the
inverse function of an algebra map $f:A\to B$. Applying the injective
map $f$ to both sides proves
\[
 g(1)=1,\qquad g(b+b')=g(b)+g(b'),\qquad
 g(bb')=g(b)g(b'),\qquad g(\iota_B(r))=\iota_A(r).
\]
For example, both sides of the multiplication equation have image $bb'$.
Thus $g$ is an algebra map, and the function inverse laws are algebra
inverse laws.
\end{proof}

\begin{proposition}[Structure identity principle]\label{prop:alguniv}
The precategory $\Alg_R$ is a category.
\end{proposition}
\begin{proof}
On a fixed set $T$, algebra structure consists of the operations
$0,1,+,\cdot,-$, a function $R\to T$, and their laws. The laws are
propositions. By universe univalence, equality of carriers is equivalent
to equivalence of carriers. Transporting the operations along a carrier
equivalence identifies algebra structures exactly when that equivalence
preserves the operations and the structure map.

This last assertion follows by equivalence induction: for the identity
equivalence, preservation says the constants and operations agree;
function extensionality makes these equalities of functions, and the
laws are propositions. The equivalence-induction argument is the
identity-elimination argument for carriers under universe univalence.

Preservation of $1,+,\cdot$ already implies preservation of $0$ and
negation, by uniqueness of additive identities and inverses. Thus the
structure-preserving equivalences are exactly algebra maps whose carrier
maps are equivalences. By Lemma~\ref{lem:algiso} they are exactly algebra
isomorphisms. This identifies the canonical function
$(A=B)\to\Iso_R(A,B)$ as an equivalence.
\end{proof}

\section{Finite presentations and a small category of algebras}

A finite presentation $p$ consists of a natural number $n$ and finitely
many polynomial equations
\[
 f_1=g_1,\ldots,f_m=g_m\qquad\text{in }R[X_1,\ldots,X_n].
\]
Write $\operatorname{Sol}_p(B)$ for the set of tuples in $B^n$ satisfying
these equations. The type $\mathcal P_R$ of such presentations belongs
to $\UU_\ell$. A presentation is \emph{realised} by an algebra $A$ if
it has a tuple $a\in\operatorname{Sol}_p(A)$ such that, for every
$R$-algebra $B$, the canonical function
\[
 \Hom_R(A,B)\longrightarrow\operatorname{Sol}_p(B),\qquad
 h\longmapsto h(a)
\]
is an equivalence. Denote the type of this data by
$\operatorname{Presents}(p,A)$, and put
\[
 \FP(A)=\tr{\sum_{p:\mathcal P_R}\operatorname{Presents}(p,A)}.
\]
This is a proposition in $\UU_{\ell+1}$: its universal property quantifies
over the algebras with small carriers.

Each $p$ has a realisation
\[
 A_p=R[X_1,\ldots,X_n]/(f_1-g_1,\ldots,f_m-g_m)
\]
with its universal tuple. Two realisations of the same presentation have
a unique isomorphism carrying one universal tuple to the other. Indeed,
the two universal properties give maps in both directions; their
composites fix the universal tuples and hence are identities.

\begin{theorem}\label{thm:fp}
There is a category $\Alg_R^{\mathrm{fp}}$ with objects and hom sets in
$\UU_\ell$ and a fully faithful inclusion into $\Alg_R$. Its object type
$I_R$ satisfies
\begin{equation}\label{eq:fpobjects}
 I_R\simeq\sum_{A:\Alg_R}\FP(A).
\end{equation}
If $a,b:I_R$ have underlying algebras $A_a,A_b$, its hom set is
$\Hom_R(A_a,A_b)$.
\end{theorem}
\begin{proof}
Apply Theorem~\ref{thm:image} to $p\mapsto A_p$. By
\eqref{eq:member}, membership of $A$ in the full image is equivalent to
\begin{equation}\label{eq:fpmember}
 \tr{\sum_{p:\mathcal P_R}\Iso_R(A_p,A)}.
\end{equation}
This proposition is equivalent to $\FP(A)$. A presentation of $A$ gives
an isomorphism $A_p\cong A$ by the common universal property. Conversely,
an isomorphism $A_p\cong A$ gives an equality of algebras by
Proposition~\ref{prop:alguniv}; transporting the known realisation of $p$
along it gives a realisation by $A$. Each use of a mere witness eliminates
into a proposition.

For completeness, the object equivalence in \eqref{eq:fpobjects} is
explicit. Send $a$ to $(ia,u_a)$, where $u_a$ is its finite-presentation proof.
Conversely, membership gives a fibre point $(a,r:ia=A)$, so send $(A,u)$
to $a$. The first round trip is the first coordinate of an equality in
a fibre of $i$, which is a proposition. The second follows from $r$ and
the fact that $\FP(A)$ is a proposition. Fullness supplies the stated
hom sets.
\end{proof}

The right side of \eqref{eq:fpobjects} is expressed in $\UU_{\ell+1}$,
but the left side belongs to $\UU_\ell$. Replacement has supplied a
small model of the type of finitely presented algebras; no propositional
resizing is used. The category contains all algebra homomorphisms
between these objects.

For example, the presentation with no generators and no relations
realises $R$, whereas the presentation with one generator $X$ and relation
$X=0$ realises $R[X]/(X)$. Evaluation at zero gives
\[
 R[X]/(X)\cong R,
\]
with inverse induced by the constants. These two codes therefore determine
equal objects in $I_R$: the isomorphism gives a path by univalence of the
small category. In general such paths retain the isomorphisms between
algebras; in particular, loops at an object retain its automorphisms.

Let $\mathcal C_R$ be the precategory with finite presentations as
objects and
\[
 \mathcal C_R(p,p')=\Hom_R(A_p,A_{p'}).
\]
Its completion is $\Alg_R^{\mathrm{fp}}$, so
Theorem~\ref{thm:descent} gives, for every category $E$,
\[
 \Fun(\Alg_R^{\mathrm{fp}},E)\simeq\Fun(\mathcal C_R,E).
\]
A functor described using presentation codes therefore descends to the
algebras themselves, with uniqueness expressed by this equivalence.
Different presentations are related through their algebra homomorphisms,
rather than by a choice of one presentation for each algebra.

\section{Set-valued functors and Yoneda}

Let $C$ be a small precategory: both its objects and its hom sets lie in
$\UU_\ell$. A covariant set-valued functor $F$ assigns a small set $F(c)$
to each object and a function $F(f):F(c)\to F(c')$ to each arrow, with
\[
 F(1_c)(x)=x,\qquad F(g\circ f)(x)=F(g)(F(f)(x)).
\]
A natural transformation $\alpha:F\to G$ consists of functions
$\alpha_c:F(c)\to G(c)$ satisfying
\[
 \alpha_{c'}(F(f)(x))=G(f)(\alpha_c(x)).
\]
Naturality is a proposition, so equality of transformations is pointwise
equality of their components. For fixed $F,G$, the set $\Nat(F,G)$ lies
in $\UU_\ell$. The type of all such functors generally lies in
$\UU_{\ell+1}$, since its values range over small sets.

For $a:C$, the covariant representable functor is
\[
 h_a(c)=C(a,c),\qquad h_a(f)(u)=f\circ u.
\]
\begin{theorem}[Yoneda]\label{thm:yoneda}
For every $a$ and every covariant set-valued functor $F$, evaluation at
the identity is an equivalence
\[
 \Nat(h_a,F)\simeq F(a),\qquad
 \alpha\longmapsto\alpha_a(1_a).
\]
Its inverse sends $x\in F(a)$ to the transformation
$u:a\to c\longmapsto F(u)(x)$.
\end{theorem}
\begin{proof}
The inverse formula is natural by the composition law for $F$.
Evaluation of this transformation at $1_a$ gives $F(1_a)(x)=x$.
Conversely, naturality of $\alpha$ at $u:a\to c$ gives
\[
 F(u)(\alpha_a(1_a))=\alpha_c(u\circ1_a)=\alpha_c(u).
\]
Pointwise equality of transformations proves the second inverse law.
\end{proof}

\begin{corollary}[Affine representables]\label{cor:affine}
On $C=\Alg_R^{\mathrm{fp}}$, let
\[
 h_a(b)=\Hom_R(A_a,A_b).
\]
Then
\[
 \Nat(h_a,h_b)\simeq\Hom_R(A_b,A_a).
\]
The inverse sends $f:A_b\to A_a$ to precomposition with $f$:
$u:A_a\to A_c\longmapsto u\circ f$.
\end{corollary}
\begin{proof}
Apply Theorem~\ref{thm:yoneda} to $F=h_b$.
\end{proof}

The representables thus give a fully faithful contravariant embedding
\[
 (\Alg_R^{\mathrm{fp}})^{\mathrm{op}}
       \longrightarrow[\Alg_R^{\mathrm{fp}},\Set_\ell].
\]
In detail, $y(1)=1$, and for $f:b\to a$ and $g:c\to b$,
\[
 y(f\circ g)=y(g)\circ y(f).
\]
These identities are respectively the unit law and associativity in the
algebra category. This is the algebraic affine functor of points: a map
out of the algebra $A_a$ into a test algebra $A_b$ is a point of $h_a$ at
$A_b$, and transformations between representables recover algebra maps
in the opposite direction. All values, natural-transformation sets, and
equivalences displayed here belong to $\UU_\ell$. This size bound uses
the small object type supplied by replacement.

Here both the represented algebras and the test algebras are finitely
presented over $R$. Full faithfulness follows from Yoneda on exactly that
category. A corresponding assertion for arbitrary schemes tested only
on finitely presented algebras requires a separate argument and suitable
hypotheses, or an enlarged test category; it does not follow from this
application of Yoneda alone.

\section{The categorical comparison for qcqs schemes}
\label{sec:qcqs}

The affine Yoneda statement above uses finitely presented test algebras.
For the scheme comparison in this section, tests are \emph{all}
commutative unital rings with set carriers in $\UU_\ell$. For such a
ring $A$, let $\operatorname{Spec}A$ be its locally ringed Zariski
lattice, and put $h_A(B)=\Hom(A,B)$. The test object type lies in
$\UU_{\ell+1}$ and its hom sets lie in $\UU_\ell$.

A geometric qcqs scheme is a small locally ringed lattice $X$ merely
admitting a finite affine cover. Its lattice and section-ring carriers
lie in $\UU_\ell$. The affine atlas is a property, rather than chosen
structure; compact overlaps need not be affine. Let $\mathcal G_\ell$
be the precategory of these objects and their local geometric maps.

For small-set-valued functors on the test rings, the Zariski lattices
form a functor $\operatorname{Zar}$. A compact open of $F$ is a natural
transformation $U:F\to\operatorname{Zar}$, with realisation
\[
 |U|(B)=\sum_{x:F(B)}(U_B(x)=\top).
\]
A functorial qcqs scheme is a Zariski sheaf merely admitting finitely
many such opens $U_i$ with pointwise join $\top$ and with each $|U_i|$
merely equivalent to some $h_{A_i}$. Sheafhood includes every finite
basic cover, including the empty cover. Let $\mathcal F_\ell$ be the
precategory of these functors and natural transformations. Its object
condition, like geometric qcqs, is a proposition.

\begin{proposition}[Univalence of both scheme categories]
\label{prop:scheme-univalence}
Both $\mathcal G_\ell$ and $\mathcal F_\ell$ are univalent categories.
\end{proposition}
\begin{proof}
For commutative rings with set carriers, and for bounded distributive
lattices, the structure identity argument of
Proposition~\ref{prop:alguniv} applies: preservation of the operations
under a carrier equivalence is the transport equation, and all laws
are propositions.

Fix a lattice $L$. An identity of ring families over $L$ corresponds
to a family of ring isomorphisms. After identifying these families,
naturality of the isomorphisms says precisely that the restriction maps
agree. Function extensionality identifies the restriction families;
their identity and composition laws are propositions. Sheafhood is
also a proposition. For a fixed presheaf, the invertibility open of a
section is unique by its defining universal property, so the type of
invertibility maps is a proposition. Their support laws are
propositions too. Thus the locally ringed conditions add no further
identity data.

A geometric isomorphism $f:X\to Y$ gives a lattice isomorphism
$f^*:L_Y\to L_X$ and compatible ring isomorphisms
\[
 f^\sharp_u:\mathcal O_Y(u)\longrightarrow
                    \mathcal O_X(f^*u).
\]
Conversely such data give an isomorphism; it and its inverse are local
because ring isomorphisms preserve and reflect units, hence preserve
the uniquely characterised invertibility opens. Apply the lattice
identity principle, retaining the inverse-image direction, then the
fixed-lattice ring-family argument above. These give contractible total
spaces of isomorphisms of locally ringed lattices. Passing between
structural isomorphisms and actual geometric isomorphisms preserves the
whole ringed map. Lemma~\ref{lem:total} therefore proves univalence
for the category of all these spaces.

For functors, an isomorphism is a natural transformation whose
components are equivalences. Pointwise univalence and function
extensionality identify its value families. After that identification,
naturality is equality of the reindexing maps; sethood and the functor
laws are propositions. Lemma~\ref{lem:total} proves univalence of the
category of small-set-valued functors.

Finally, a full subcategory defined by a propositional object condition
inherits univalence. Its isomorphisms are the original ones, and proofs
of the object condition add no identity data. Apply this to the two
qcqs conditions to obtain the assertion.
\end{proof}

The geometric category has objects in $\UU_{\ell+1}$ and hom sets in
$\UU_\ell$. For the functorial category both bounds are initially
$\UU_{\ell+1}$. We may therefore place both categories at the common
bound $\ell+1$; the comparison uses no resizing of natural-transformation
types.

\begin{theorem}[Equivalence of qcqs scheme categories]
\label{thm:qcqs-comparison}
The functor of points
\[
 H:\mathcal G_\ell\longrightarrow\mathcal F_\ell,\qquad
 h_X(B)=\Hom(\operatorname{Spec}B,X),\qquad
 H(f)_B(p)=f\circ p,
\]
is an equivalence of univalent categories. In particular,
\[
 \Ob(\mathcal G_\ell)\simeq\Ob(\mathcal F_\ell).
\]
It has an explicit inverse functor and natural isomorphisms for both
composites, without an axiom of choice.
\end{theorem}
\begin{proof}
The geometric sheaf and atlas theorems make $h_X$ a functorial qcqs
scheme. The geometric unit and associativity laws make $H$ a functor.
The checked absolute Hom--Nat theorem gives full faithfulness:
\[
 \Hom(X,Y)\simeq\Nat(h_X,h_Y).
\]
The checked absolute essential-image theorem gives, for each
$F\in\mathcal F_\ell$,
\[
 \tr{\sum_{X:\mathcal G_\ell}(F\simeq h_X)}.
\]
Here the displayed equivalences are natural. Inverting them gives mere
essential surjectivity of $H$.

For clarity, this input reconstructs an actual finite atlas using
geometric transition isomorphisms on compact overlaps. Their triple
domains are identified before the cocycle equations are compared.
Finite gluing gives a small geometric qcqs space, and compatibility in
both directions identifies its functor with $F$. The merely given
atlas and affine witnesses are eliminated only into mere
representability. The companion note, \emph{The essential image of the
functor of points}, gives this construction and its detailed proof.

Proposition~\ref{prop:scheme-univalence} supplies univalence of both
categories. Theorem~\ref{thm:weak-equivalence} now applies to $H$.
In particular it makes each representation fibre a proposition before
removing the truncation. Formula~\eqref{eq:inverse-arrow} gives the
inverse on arrows, with its functor laws and both natural isomorphisms.
\end{proof}

The theorem retains automorphisms as loops in the object types. It
compares geometric schemes with functorial schemes, rather than passing
to their sets of isomorphism classes. The theorem also has the following
relative form, obtained by totalising the native functor on algebra tests.

\begin{theorem}[Relative qcqs comparison]\label{thm:relative-comparison}
Let $k$ be a small commutative unital ring with set carrier. Let
$\mathcal F_{k,\ell}$ be the category of native small-set-valued functors
on all small $k$-algebras satisfying the relative Zariski sheaf and finite
affine compact-open atlas conditions. The relative functor of points
\[
 H_k:\mathcal G_\ell/\operatorname{Spec}k
             \longrightarrow\mathcal F_{k,\ell},\qquad
 H_k(X,s)(B,\iota)=\Hom_{\operatorname{Spec}k}
           (\operatorname{Spec}(B,\iota),(X,s)),
\]
is an equivalence of univalent categories, with an explicit inverse
functor and both natural isomorphisms.
\end{theorem}
\begin{proof}
For a native relative functor $F$, define its totalisation by
\[
 (\mathsf T F)(B)=\sum_{\iota:k\to B}F(B,\iota).
\]
On $h:B\to C$, reindex $(\iota,x)$ along the canonical algebra map
$c_{h,\iota}:(B,\iota)\to(C,h\iota)$. If $p:h\iota=\kappa$, path
induction gives the precise dependent equation
\[
 \operatorname{transport}_{j\mapsto F(C,j)}
       (p,F(c_{h,\iota})(x))=F(h,p)(x).
\]
Thus total reindexing on any actual algebra map $a$ identifies the whole
point with $(\kappa,F(a)(x))$. This proves the total functor laws and
naturality of its projection $\pi_F:\mathsf T F\to h_k$.

Intrinsic relative sheafhood is equivalent to absolute sheafhood of
$\mathsf T F$. For the forward direction, project a compatible total
cover family $(\iota_i,x_i)$ to base maps and glue them in the sheaf
$h_k$. Rebase each $x_i$ along the resulting equation
$\lambda_i\iota=\iota_i$ before applying the relative sheaf condition.
The whole-pair reindexing equation gives compatibility on the fixed
base, and base uniqueness followed by relative uniqueness proves total
uniqueness. Conversely, total gluing of paired relative values and base
separation in $h_k$ give a value in the specified fibre. Injectivity of
pairing at that fixed base proves relative uniqueness. Empty covers
are included in both arguments.

Compact opens correspond by evaluation at $(\iota,x)$, and realisations
commute with totalisation by reordering the membership proof. Totalising
a relative affine representable gives $h_A$: its ring map $h:A\to B$
forces the target base map to be $h\alpha$, where $\alpha:k\to A$ is
the source structure. Conversely, an absolute affine chart gives
$\alpha$ by projecting its generic point. Naturality supplies the base
path needed to transport its values into each specified relative
fibre. Thus relative and total affine atlas conditions agree. Together,
these prove
\[
 \text{$F$ is relative functorially qcqs}
 \quad\Longleftrightarrow\quad
 \text{$\mathsf T F$ is absolute functorially qcqs}.
\]

Apply the absolute essential-image theorem to $\mathsf T F$. For an
actual representing equivalence $e:\mathsf T F\simeq h_X$, the natural
projection $\pi_F e^{-1}:h_X\to h_k$ gives $s:X\to\operatorname{Spec}k$
by the affine comparison and full faithfulness. Its equation on affine
probes identifies the projection fibres with points over $k$. On the
original side the fibre over $\iota$ is
\[
 \sum_{\kappa:k\to B}\sum_{x:F(B,\kappa)}(\kappa=\iota)
       \simeq F(B,\iota),
\]
by transport along the displayed path. Its inverse pairs with $\iota$
and reflexivity; path induction proves both inverse laws. Restricting
$e$ to these fibres gives a native relative representing equivalence.
Only mere relative representability is asserted before eliminating the
mere absolute representation. Conversely total relative geometric
points are $h_X$, since each ordinary point has a unique base lift.
The totalisation criterion gives the relative geometric qcqs property.
This proves mere essential surjectivity of $H_k$. The existing relative
Hom--Nat theorem gives full faithfulness on the same native functors.

For univalence, the geometric category is the slice of a univalent
category. Its isomorphisms are base isomorphisms $e:x\cong y$ satisfying
$te=s$; their inverses commute since $se^{-1}=t$. Over the identity
base isomorphism, the compatible structure arrows form the contractible
singleton $(t,t1=s)$. The total-space criterion therefore proves slice
univalence. On the functorial side, the value-family and action identity
argument applies to the fixed type of actual algebra tests, and the
relative qcqs condition is a proposition. Both categories are univalent.
Theorem~\ref{thm:weak-equivalence} now supplies the stated equivalence
and its explicit quasi-inverse.
\end{proof}

The detailed transport, rebasing, and fibre arguments, with their
correspondence to the checked declarations, are given in the companion
note \emph{Relative totalisation over a base ring}. The native sheaf
and atlas conditions are defined on algebra tests before the
intrinsic-total equivalence is proved; they are not abbreviations for
absolute conditions on $\mathsf T F$.

\section{Connection with the implementation}\label{sec:implementation}

The preceding constructions are implemented in Unit, in
\code{sag/categories/}, \code{sag/fpalg/}, and \code{sag/qcqs/}.
The categorical foundations check with 21 modules, and the finitely
presented algebra development with 47 modules, including dependencies.
The complete absolute and relative qcqs comparison checks with
249 modules and 9,250 declarations.
All three checkpoints have zero failures. The informal arguments are
recorded in the design documents and at the heads of the source files.
The correspondence is as follows.

\medskip
\noindent\begin{tabular}{@{}>{\raggedright\arraybackslash}p{.39\textwidth}%
>{\raggedright\arraybackslash}p{.57\textwidth}@{}}
\textbf{Mathematical construction}&\textbf{Source files}\\[4pt]
Categories and the total-space criterion&
\code{Category.seg}\\[4pt]
Functors and natural transformations&
\code{Functor.seg}, \code{FunctorCat.seg}\\[4pt]
Small full image (Theorem~\ref{thm:image})&
\code{SmallImage.seg}\\[4pt]
Full faithfulness and descent&
\code{PrecomposeFF.seg}, \code{Descent.seg},
\code{UniversalProperty.seg}\\[4pt]
Univalence and the object equivalence&
\code{FunctorUnivalence.seg}, \code{EquivObjects.seg}\\[4pt]
Completion property \eqref{eq:completion}&
\code{CompletionUP.seg}\\[4pt]
Algebra category and structure identity&
\code{AlgebraCategory.seg}, \code{AlgebraSIP.seg}\\[4pt]
Small finitely presented algebra category&
\code{FPImage.seg}, \code{FPComparison.seg}\\[4pt]
Descent from presentation codes&\code{FPDescent.seg}\\[4pt]
Yoneda and affine representables&
\code{Yoneda.seg}, \code{FPYoneda.seg}\\[4pt]
Equivalence from mere essential surjectivity&
\code{WeakCategoryEquivalence.seg}, \code{EmbeddingFibres.seg}\\[4pt]
Explicit inverse and natural isomorphisms&
\code{CategoryQuasiInverse.seg}\\[4pt]
Univalent geometric and functorial categories&
\code{GeometricCategory.seg}, \code{FunctorCategory.seg},
\code{FullSubcategory.seg}\\[4pt]
Qcqs comparison (Theorem~\ref{thm:qcqs-comparison})&
\code{QcqsCategories.seg}
\end{tabular}
\medskip

The category interface is \code{Category[k,l]}: objects have level $k$
and hom sets level $l$. For a family indexed at level $l$ in such a
category, \code{SmallImage.SmallFullImage} gives a
\code{Category[l,l]}. Its image fibre is a proposition, and
\code{from-fibre} and \code{to-fibre} express \eqref{eq:member}.
The completion theorem exposes both full faithfulness and an equivalence
on functor types. Functor categories are placed in a common universe
bounding the source and target levels.

For algebras, \code{AlgebraSIP.Over.category} has levels $l+1,l$.
The algebra representation and presentation universal properties are
reused from \code{sag/presentations/}; the structure identity proof uses
\code{SipGeneric} and local structure pieces for negation and the map
from the base ring. \code{FPImage.Over.SFI.category} has levels $l,l$.
The statements \code{FPComparison.Over.comparison} and
\code{FPComparison.Over.objects} give respectively membership
$\simeq\FP(A)$ and \eqref{eq:fpobjects}. Thus the comparison is with the
existing universal-property definition of finite presentation, not merely
with a new predicate on codes.

Set-valued functors are represented directly by their values, functions
on arrows, and pointwise laws. Natural transformations are component
families with pointwise naturality; \code{nat-eq} proves their
extensionality. This interface suffices for \code{Yoneda.yoneda}, and
\code{FPYoneda} states the instance in algebra-homomorphism notation.
The precomposition map is proved to be an equivalence, and its identity
and reversed composition laws are checked. No packaging of all these
set-valued functors as a univalent category of functors into a category
of sets is needed for these results.

The foundational parameters are universe univalence and the two parts
of replacement: small images and their surjectivity elimination
principle. Function extensionality and propositional univalence are
derived. The small full image uses replacement once; the algebra
application uses it at levels $l,l+1$. The object equivalence does not
use propositional resizing. Later Yoneda arguments use the small category
already constructed and require no further completion.

\begingroup\raggedright
For the scheme comparison, the sources in \code{sag/qcqs/} reuse
\code{SipGeneric} and the existing category interfaces through local
symlinks, without modifying their implementations.
\code{RingIdentity} and \code{LatticeIdentity} prove the algebraic
identity principles. \code{FamilyIdentity} and \code{DecoratedIdentity}
pass them to indexed families, structure, and propositional properties.
\code{PresheafIdentity} and \code{LocalIdentity} handle restrictions,
sheafhood, and locality. \code{SpaceIdentity} and
\code{SpaceIsoIdentity} connect the structural presentation to the
actual geometric maps, with an equation of their underlying ringed maps.
The total-space criterion then gives \code{GeometricCategory.univalent}.
\code{SetIdentity} and \code{FunctorIdentity} supply the corresponding
argument for \code{FunctorCategory.univalent}.

\code{QcqsCategories.functor} is the existing functor of points;
\code{fully-faithful} uses \code{PointsUniversal}, and
\code{essentially-surjective} uses \code{EssentialImage.reconstruct}.
The source and target are actual univalent category records, obtained
by \code{FullSubcategory}. Their univalence proofs are constructed,
rather than additional parameters of the comparison.
\code{QcqsCategories.object-equivalence} records the equivalence of
object types. \code{QcqsCategories.equivalence} packages the explicit
inverse functor, its identity and composition laws, and the two natural
isomorphisms, including both naturality proofs.

For Theorem~\ref{thm:relative-comparison},
\code{RelativeTotal.change-target} states the displayed transport
comparison and \code{on-algebra} identifies the whole total point on
an actual algebra map. \code{RelativeBaseGlue},
\code{RelativeTotalSheaf}, and \code{RelativeSheafOfTotal} implement
base gluing, rebasing and the two sheaf implications.
\code{RelativeTotalisation} exposes the sheaf and qcqs equivalences;
\code{RelativeAffineOfTotal} recovers the generic algebra structure.
\code{RelativeFibreRecovery}, \code{RelativeFibreEquivalence}, and
\code{RelativeGeometricFibre} prove the fibre identifications, retaining
the dependent base equations. \code{RelativeEssentialImage} proves
both directions of relative representation.

\code{SliceCategory} proves univalence of the geometric slice, and
\code{RelativeFunctorIdentity} with \code{RelativeFunctorCategory}
proves native functor-category univalence. The final statement is
\code{RelativeCategories.Over[K].equivalence}, with the inverse functor
and both natural isomorphisms supplied by the same general
\code{CategoryQuasiInverse} used absolutely.

The full independent qcqs check has 9,250 declarations and zero failures.
The combined \code{essential-image.opens}, \code{categories.opens}, and
\code{relative.opens} audit has 3,752 opening lists; every reported list
says that nothing else opened. The absolute categorical continuation
has 591 declarations in 18 modules and reuses 167 declarations in five
source modules. The relative continuation adds 863 audited declarations
in 35 modules. No core or category refactoring, new axiom, choice
principle, or resizing principle is used.

Thus this checkpoint includes the categorical foundations and both the
absolute and relative equivalences of geometric and functorial qcqs
schemes.
\par\endgroup

\end{document}
